\begin{abstract}
A procedure can be correct where it was learned yet incomplete where it is reused. We study coding agents supplied with procedural memory in target tasks that change a security relevant source assumption. Thirteen synthetic families cross four memory conditions and two repetitions in six model and agent configurations, yielding \RecordedRuns{} recorded runs and \ValidRuns{} technically valid outcomes. The primary endpoint is functional completion that fails a focal security witness. The effect of correct memory ranges from $\PrimaryURangeLow$ to $\PrimaryURangeHigh$ percentage points across configurations, with positive, zero, and negative estimates. Lower endpoint values sometimes accompany lower functionality. A generic applicability reminder reduces the endpoint in all three Codex configurations, a configuration specific direction that survives every assignment of unresolved outcomes in the recorded cohort; the MiniSWE configurations do not share that pattern. In a prespecified 52-run sample, two model assisted coding passes agree on 22 implementations with a target protection but on only seven with explicit assumption recognition in the visible trace before editing. Matched transcripts and patches show unchecked reuse, successful adaptation, and distinct mechanisms behind the same measured failure. Correct procedural memory is therefore neither uniformly beneficial nor uniformly harmful when its source assumptions change, and outcome evaluation must observe functionality and the focal property separately.
\end{abstract}
